\section{Exact quadratic Schur tests of the residual comparison pole}
\label{sec:s5-quadratic-schur}

Retain the actual nontrivial branch of MA-RB-002 at its real
rightmost pole \(\gamma\in(1/2,1)\), with pole order \(k\ge1\).
Write
\[
 \begin{gathered}
 b=\operatorname{ord}_\gamma C,\quad
 n=\operatorname{ord}_\gamma A,\quad
 v_0=\operatorname{ord}_\gamma L_V,\quad
 v_1=\operatorname{ord}_\gamma L_{V\eta},\quad
 e=\operatorname{ord}_\gamma L_X,\\
 x=\operatorname{ord}_\gamma L_W=n+k,\qquad
 x_1=\operatorname{ord}_\gamma L_{W\eta}.
 \end{gathered}
 \tag{MA-RB.46}
\]
All these orders are nonnegative. The justification for \(x\)
is the actual comparison pole, not a general assumption of
Artin holomorphy. At this real point
\(L_{W\eta}=\zeta_F/\zeta_{\mathbb Q}\) is holomorphic, since
\(\zeta_{\mathbb Q}(\gamma)\ne0\).

Let \(\Pi,\nu,\varpi\) be Calegari's four-dimensional
representation, its finite-order similitude character, and its
specified five-dimensional comparison. A twist written outside
\(\operatorname{Sym}^2\) means a twist of that \(L\)-function.
In particular \(\operatorname{Sym}^2(\varpi)\otimes\eta\)
denotes a different operation from
\(\operatorname{Sym}^2(\varpi\otimes\eta)\).

\paragraph{MA-RB-005 (the quadratic Schur holomorphy constraints).}
Choose a finite set \(S\) containing infinity, \(2\), and every
ramified or omitted place in the comparisons. Define
\[
 \begin{aligned}
 Q_4(s)&=L^S(s,\Pi,\operatorname{Sym}^2\otimes\nu^{-1}),&
 Q_{4,\eta}(s)&=L^S(s,\Pi,\operatorname{Sym}^2
                                      \otimes\nu^{-1}\eta),\\
 Q_5(s)&=L^S(s,\varpi,\operatorname{Sym}^2),&
 Q_{5,\eta}(s)&=L^S(s,\varpi,\operatorname{Sym}^2\otimes\eta).
 \end{aligned}
 \tag{MA-RB.47}
\]
All four functions are holomorphic at \(\gamma\), and
\[
 \begin{array}{c|c}
 \text{function}&\text{order at }\gamma\\ \hline
 Q_4&v_1+e-k\\
 Q_{4,\eta}&v_0+e+k\\
 Q_5&v_0+n+x_1\\
 Q_{5,\eta}&b+v_1+n+x_1+2k .
 \end{array}
 \tag{MA-RB.48}
\]
The partial Euler products for \(\wedge^2\varpi\)
and \((\wedge^2\varpi)\otimes\eta\) are \(Q_4\) and
\(Q_{4,\eta}\). This supplies their continuation and
holomorphy without positing a separate automorphic
exterior-square lift of \(\varpi\).

Among the four order nonnegativities, only \(v_1+e\ge k\)
is not automatic from (MA-RB.46); that is already MA-RB-002.
Thus these quadratic constraints do not increase the
coefficient nine in its normal-closure multiplicity bound
or exclude the formal genuine-character assignment of
MA-RB-004. They do force the additional actual zeros
\[
 \operatorname{ord}_\gamma Q_{4,\eta}\ge k,\qquad
 \operatorname{ord}_\gamma Q_{5,\eta}\ge2k .
 \tag{MA-RB.49}
\]
No nonvanishing of either function is asserted.

\begin{proof}
\emph{Finite representation identities.}
Use the already checked \(S_5\) character table, with columns
\(\mathbf1,\eta,V,V\eta,W,W\eta,X\). The relevant new
calculations, in the specified class order, are
\[
 \begin{array}{c|rrrrrrr}
 &1&2A&2B&3A&6A&4A&5A\\ \hline
 \chi_W&5&1&1&-1&1&-1&0\\
 \chi_W(g^2)&5&5&5&-1&-1&1&0\\
 \chi_{\operatorname{Sym}^2W}&15&3&3&0&0&1&0\\
 \chi_{\wedge^2W}&10&-2&-2&1&1&0&0 .
 \end{array}
 \tag{MA-RB.50}
\]
For eigenvalues \(\lambda_j\), the sums of
\(\lambda_j\lambda_l\) over \(j\le l\) and \(j<l\)
are \((\chi_W(g)^2+\chi_W(g^2))/2\) and
\((\chi_W(g)^2-\chi_W(g^2))/2\), respectively. They
give every entry displayed. Adding columns of the
checked character table now proves
\[
 \operatorname{Sym}^2W
   =\mathbf1\oplus V\oplus W\oplus W\eta,\qquad
 \wedge^2W=V\eta\oplus X .
 \tag{MA-RB.51}
\]
Equality of characters proves the representation identities.
The table also gives \(X\eta=X\), so their external
\(\eta\)-twists have constituents
\(\eta,V\eta,W,W\eta\) and \(V,X\), respectively.
This reuses the existing table without reopening the
earlier finite representation certificates.

\emph{The general symplectic parameter identity.}
Calegari's Lemma~4.2 and the subsequent parameter formula
give, at every unramified place,
\[
 S_4=\{a,b,\nu/a,\nu/b\},\qquad
 T_5=\{1,ab/\nu,\nu/(ab),a/b,b/a\}.
 \tag{MA-RB.52}
\]
The eigenvalues of
\(\nu^{-1}\operatorname{Sym}^2 S_4\) are
\[
 \{a^2/\nu,b^2/\nu,\nu/a^2,\nu/b^2,
                   ab/\nu,a/b,b/a,\nu/(ab),1,1\}.
 \tag{MA-RB.53}
\]
The products of two distinct entries of \(T_5\) give
the same multiset: the pairs containing \(1\) give the
four middle entries, the two inverse pairs give \(1,1\),
and the remaining four pairs give the first four entries.
Hence
\[
                  \wedge^2T_5=\nu^{-1}\operatorname{Sym}^2S_4 .
 \tag{MA-RB.54}
\]
Multiplication by \(\eta(p)\) gives the external twist.
Taking reciprocal Euler polynomials outside \(S\) proves
the asserted identities with \(Q_4,Q_{4,\eta}\) initially
in an absolute-convergence half-plane. Equality with
these known functions then supplies the meromorphic
continuations of the exterior-square products. No
additional functorial lift is assumed.

\emph{Exceptional Euler factors.}
At \(p\in\Omega\) the class is \(4A\), and the two
five-dimensional multisets are
\[
        T_A=\{1,1,1,i,-i\},\qquad
        T_W=\{1,-1,-1,i,-i\}.
 \tag{MA-RB.55}
\]
For \(P_{\mathcal E}(T)=\prod_{\lambda\in\mathcal E}
(1-\lambda T)\), direct multiplication gives
\[
 \begin{array}{c|c|c}
 &\text{automorphic polynomial}&\text{Artin polynomial}\\ \hline
 \operatorname{Sym}^2&
 (1-T)^7(1+T)^2(1+T^2)^3&
 (1-T)^5(1+T)^4(1+T^2)^3\\
 \wedge^2&
 (1-T)^4(1+T^2)^3&
 (1-T)^2(1+T)^2(1+T^2)^3 .
 \end{array}
 \tag{MA-RB.56}
\]
Indeed, the multiplicities of \(1,-1,i,-i\) in the
two symmetric squares are \(7,2,3,3\) and \(5,4,3,3\).
In the exterior squares they are \(4,0,3,3\) and
\(2,2,3,3\). These counts prove all signs in (MA-RB.56).
The \(L\)-factor is \(P^{-1}\), so its automorphic to
Artin ratio in either row is
\[
       \Delta(T)=\frac{P_{\rm Artin}(T)}{P_{\rm automorphic}(T)}
                   =\left(\frac{1+T}{1-T}\right)^2 .
 \tag{MA-RB.57}
\]
Since \(\eta(p)=-1\) on \(4A\), the external twist
replaces \(T\) by \(-T\) and the ratio by \(\Delta^{-1}\).
At all other unramified places the original parameters
agree; so do their Schur functors and external twists.

Put \(D^S(s)=\prod_{p\in\Omega\setminus S}\Delta(p^{-s})\).
Removing finitely many primes of \(\Omega\) changes
\(D_\Omega\) by a holomorphic nonzero factor at
\(\gamma>0\); thus \(\operatorname{ord}_\gamma D^S=-k\).
Combining (MA-RB.51)--(MA-RB.57) gives
\[
 \begin{aligned}
 Q_4&=D^S L_{V\eta}^S L_X^S,&
 Q_{4,\eta}&=(D^S)^{-1}L_V^S L_X^S,\\
 Q_5&=D^S\zeta_{\mathbb Q}^S L_V^S L_W^S L_{W\eta}^S,&
 Q_{5,\eta}&=(D^S)^{-1}C^S L_{V\eta}^S L_W^S L_{W\eta}^S .
 \end{aligned}
 \tag{MA-RB.58}
\]
These are first identities of absolutely convergent Euler
products and then meromorphic identities. All removed
finite Artin factors are units at \(\gamma\): their
inverse Euler polynomials have roots of modulus at least
one, whereas \(p^{-\gamma}<1\). We use the four
\emph{partial} symmetric-square functions; no assertion
about unspecified ramified factors of a new lift is needed.
Deletion of \(S\) therefore changes none of the Artin orders
on the right. Also \(\operatorname{ord}_\gamma
\zeta_{\mathbb Q}^S=0\).

\emph{Known holomorphy and the exact orders.}
Takeda's Theorem~7.1 applies to both cuspidal representations
\(\Pi,\varpi\), with the four unitary twists in (MA-RB.47).
Their central characters are unitary and \(\nu,\eta\)
have finite order. These are the already checked
Calegari inputs: the finite-image base changes have
finite-order central character, hence the descended
central characters are unitary. Enlarge \(S\) as in
Takeda's theorem. Its only possible poles occur at
zero and one, so all four functions are holomorphic at
\(\gamma\in(1/2,1)\).
Taking orders in (MA-RB.58) gives
\[
 \begin{split}
 \operatorname{ord}_\gamma Q_4&=v_1+e-k,\\
 \operatorname{ord}_\gamma Q_{4,\eta}&=v_0+e+k,\\
 \operatorname{ord}_\gamma Q_5&=v_0+x+x_1-k=v_0+n+x_1,\\
 \operatorname{ord}_\gamma Q_{5,\eta}
             &=b+v_1+x+x_1+k=b+v_1+n+x_1+2k .
 \end{split}
 \tag{MA-RB.59}
\]
This proves (MA-RB.48)--(MA-RB.49). All functions used
here are genuine meromorphic \(L\)-functions or exact
Euler products identified with them; these are actual
pointwise identities, not freely assigned analytic orders.

\emph{Scope of these inequalities.}
The last three expressions in (MA-RB.59) are
automatically nonnegative under (MA-RB.46). The first
gives exactly \(v_1+e\ge k\). For every positive integer
\(k\), the formal assignment of MA-RB-004 is
\[
       (b,n,v_0,v_1,e,x,x_1)=(0,0,0,k,0,k,0).
\]
It gives
\[
 (\operatorname{ord}Q_4,\operatorname{ord}Q_{4,\eta},
       \operatorname{ord}Q_5,\operatorname{ord}Q_{5,\eta})
      =(0,k,0,3k),\qquad
 b+4v_0+4v_1+5x+5x_1+6e=9k .
 \tag{MA-RB.60}
\]
Thus it satisfies all the new inequalities and attains
the old regular-order lower bound. This proves the
specified failure to obtain a stricter bound or a
contradiction from this family of nonnegativities.
It does not assert arithmetic realization of that
formal assignment.
\end{proof}

\paragraph{Higher exterior powers of the same parameters.}
The multiset \(T_5\) in (MA-RB.52) is invariant under
inversion and has product one. Complementing a subset
therefore gives \(\wedge^{5-r}T_5=\wedge^rT_5\).
The same holds for \(W\): it is self-dual and its
determinant is one, since a transposition has three
eigenvalues \(1\) and two eigenvalues \(-1\), and
transpositions generate \(S_5\). Thus its third exterior
power repeats the second, its fourth repeats the
standard representation, and its fifth is trivial.
For \(S_4\), its product is \(\nu^2\) and its inverse
multiset is \(\nu^{-1}S_4\), so
\(\wedge^3S_4=\nu S_4\) and \(\wedge^4S_4=\nu^2\).
These introduce no further standard exterior-power
factor beyond the preceding ones and the original
standard functions. No theorem about higher symmetric
powers is inferred from Takeda's quadratic theorem.

\paragraph{Primary inputs and actual verification scope.}
The new character calculation uses Calegari's character
table in \S2. His Lemma~4.2 and the following displayed
decomposition and parameter formula supply (MA-RB.52);
Lemma~4.6 gives \(\nu=\psi^{-1}\); Theorem~4.7 and the
local calculation in Lemma~4.8 give (MA-RB.55).
The supplied primary LaTeX locators are
\texttt{lemma:satake}, \texttt{lemma:sim}, the theorem
immediately following the latter, and \texttt{lemma:aa}.
The exact source spans read for this deduction are
lines 554--602 and 735--807 of \texttt{S5ArtinWeb.tex}.
These locators do not conflate article-PDF and journal pages.
The literal base-change summand error newly detected in
the proof of Lemma~4.2 is recorded and repaired in
MA-S5-008. That complete repair preserves the symplectic
parameter input used here; the printed identical-summand
argument is not being endorsed.

The analytic input is Takeda (2015), Theorem~7.1,
\texttt{arXiv:1506.04791v1}, printed pp.~49--50.
Its entire terminal proof was read in the MA-RB-002 audit.
Here its hypotheses are checked for the additional
five-dimensional representation and the external twists.
The upstream Eisenstein-series and integral-representation
results remain the declared classical inputs; no new
full audit of those results is claimed.

The exact new calculations are reproducible in
\texttt{checks/s5/s5\_quadratic\_schur.py} and its JSON
output. They check the two character decompositions,
the symbolic \(\operatorname{GSp}_4\) weight identity,
the four exceptional ratios and the order expressions,
using the previously checked finite table.
The written proof is independent of floating-point
calculations. In the Artin-motive language fixed in
MA-S5, (MA-RB.51) is an identity of symmetric and exterior
operations in the representation \(\lambda\)-ring.
The inherited exact locators remain Deligne (1979),
Definition~6.1 and \S\S6.2--6.3, and SGA~1, Expos\'e~V,
\S7. This interpretation supplies no missing
nonvanishing at \(\gamma\).

Calegari, F. (2013). The Artin conjecture for some
\(S_5\)-extensions. \emph{Mathematische Annalen, 356},
191--207. \url{https://doi.org/10.1007/s00208-012-0839-4}.

Takeda, S. (2015). On a certain metaplectic Eisenstein series
and the twisted symmetric square \(L\)-function.
\emph{Mathematische Zeitschrift, 281}, 103--157.
\url{https://doi.org/10.1007/s00209-015-1476-x}.
